\section{Discussion}\label{sec:discussion}

By combining nonparametric partial identification, composition-adjusted missingness diagnostics, 30-hypothesis FWER inference and cost-conditional follow-up benchmarks, this study delineates what CDC NORS foodborne outbreak records establish robustly from what depends on missing-data assumptions, case ascertainment and finite-sample error. Without assumptions on unrecorded hospitalization outcomes, 4 of 15 sample orderings separate and STEC above norovirus GII remains significant across all 30 two-sided tests under every inferential check ($p = 0.024$ outbreak-HC1; $p = 0.031$ bootstrap covariance; $p = 0.033$ 53-state clustering; $p = 0.031$ wild cluster bootstrap; $p = 0.043$ non-Gaussian studentized bootstrap). High stratum concentration ($\text{HHI} \approx 0.50$, $N_{\text{eff}} \approx 2.0$ in sparse strata) inflates upper-bound variance for \textit{C. perfringens} and norovirus GI, leaving simultaneous worst-case ranks at $[1, 5]$ for STEC, $[2, 6]$ for norovirus GII and $[1, 6]$ for the remaining pathogens.

Commonly reported denominator conventions map onto specific slices of our framework: dividing hospitalizations by all ill persons ($H/C$) matches worst-case lower bounds $L_j$~\cite{cdc2026surveillance2019,heiman2015escherichia}, whereas dividing by persons with known status ($H/K$) approximates the within-cell missing-at-random anchor ($\delta = 0$)~\cite{burke2019norovirus,wikswo2022enteric}. In comparison with US national burden estimates by Scallan et al.~\cite{scallan2011foodborne} (updated for 2019 in~\cite{scallanwalter2025foodborne}) and prior NORS summaries~\cite{deweymattia2018surveillance}, our standardized $\delta = 0$ proportions for STEC (28.5\%), \textit{Salmonella} (17.6\%) and \textit{Campylobacter} (8.5\%) fall below their proportions among laboratory-confirmed bacterial illnesses (46.2\% for O157 STEC, 27.2\% for \textit{Salmonella} and 17.1\% for \textit{Campylobacter}, reflecting broader non-O157 STEC inclusion and NORS setting standardization), while matching \textit{C. perfringens} (0.6\% vs.\ 0.6\% derived from outbreak data). For norovirus, however, Scallan et al.\ estimate only 0.03\% hospitalizations per total estimated community illness, whereas NORS outbreak-reported proportions are roughly 50-fold higher (1.6\% for norovirus GII at $\delta = 0$): milder viral gastroenteritis outbreaks seldom reach public health surveillance unless severe cases trigger investigation. Across all three lowest-risk agents, worst-case ambiguity (15.9--19.3 points) exceeds 1 percentage-point differences, rendering their internal ordering unresolvable ($\delta^{*} = 0.19$--$4.2$ points)~\cite{vega2014genotypic,grass2013epidemiology,burke2019norovirus}. Discarding the 157 incomplete records (variant R3b) would overstate sample worst-case orderings (7 vs.\ 4) by changing the estimand to omit ambiguous records.

Crucially, stratified checks clarify why crude comparisons between partially and fully recorded outbreaks must be interpreted cautiously. Although crude known-status hospitalization was higher in partially recorded outbreaks for STEC and \textit{Salmonella} and tracked rates in the 11 low-missingness states to fully recorded levels (26.2\% STEC; 13.0\% \textit{Salmonella}), this gap was largely confounded by outbreak size and multistate laboratory surveillance: within cells, the STEC gap vanished ($+0.19$ points; ``Std obs $-$ Full'' $= -3.22$ points). Thus, the data do not establish a uniform direction of residual missingness departure. Nevertheless, all 12 upper-tier orderings remain robustly separated across $\pm 5$-point symmetric bands, $-15/+5$ to $-20/+5$ asymmetric stress tests, and odds-ratio factors up to $\Gamma = 3.0$ (14 of 15 total pairs including two lower-tier pairs at $\Gamma = 3.0$; 12 upper-tier pairs throughout $\Gamma \le 5.0$). Simultaneous 95\% rank sets computed at worst case, $\delta = 5$ points and $-15/+5$ points identically separate STEC $[1, 1]$, \textit{Salmonella} $[2, 2]$ and \textit{Campylobacter} $[3, 3]$ from the lower tier $[4, 6]$ whenever departure bounds hold at those scenarios.

Operational follow-up triage depends critically on how agencies incur verification costs. Targeting follow-up by stratum leverage $w_s/C_{js}$ dominates when costs scale per case resolved, whereas targeting by outbreak width $w_s U_i/C_{js}$ dominates when contacting a jurisdiction carries a fixed per-outbreak cost (under which targeting largest outbreaks also beats random selection). Framing prioritization around expected bound reduction per unit cost resolves this apparent reversal transparently across all cost regimes.

\subsection{Limitations}\label{sec:limitations}

First, NORS reporting is voluntary and overall outbreak reporting rates vary four-fold between high- and low-reporting states (spanning roughly 5 to 86 outbreaks per 10 million population annually, and $>40$-fold for norovirus)~\cite{white2022foodborne}. Second, partial-identification bounds take estimated case counts $C_i$ (as well as $H_i$ and $K_i$) as accurate; differential case ascertainment between laboratory-confirmed multistate clusters (common in STEC and \textit{Salmonella}) and complaint-detected local outbreaks can bias rates by lowering detected mild illnesses~\cite{heiman2015escherichia}, though restricting entirely to single-state outbreaks (variant R11) preserved every main verdict. Third, common support was selected data-dependently ($\ge 3$ outbreaks per group), dropping 30.0\% of cases and placing 90.4\% of pooled weight in 2009--2018 because only restaurants met support criteria in 2019--2023 (9.6\%); retrospective follow-up of 2009--2018 records is practically infeasible, so our follow-up simulation applies prospectively to ongoing surveillance. Fourth, the masking cohort of fully recorded outbreaks is smaller and more often single-state than partially recorded targets. Fifth, outbreak-versus-sporadic hospitalization comparability in Ebel et al.~\cite{ebel2016comparing} pertains strictly to laboratory-confirmed FoodNet bacterial infections (\textit{Salmonella}, O157 STEC, \textit{Campylobacter}, \textit{Listeria}) and not non-O157 STEC, \textit{C. perfringens} or norovirus; our genus-level estimand lacks age/comorbidity covariates and pools distinct serotypes and subtypes~\cite{jones2008salmonellosis,gould2013increased,geissler2017increasing}. Sixth, bootstrap-world calibration showed anti-conservative 30-test rejection rates for asymptotic HC1 in sparse strata ($N_{\text{eff}} \approx 2$), and wildly unequal state cluster sizes (3 states holding 91.4\% of unrecorded STEC cases) can cause wild cluster bootstrap tests to slightly over-reject~\cite{mackinnon2017wild}.

\subsection{Conclusions}\label{sec:conclusions}

In US foodborne outbreak surveillance (2009--2023, predominantly weighted to 2009--2018), STEC exhibits a higher standardized hospitalization proportion than norovirus GII without unrecorded-outcome assumptions across all 30 two-sided tests and five inferential procedures. Under $\delta = 5$ percentage points or asymmetric $-15/+5$ percentage-point stress tests, all 12 upper-tier orderings hold jointly and yield simultaneous 95\% rank sets separating STEC $[1,1]$, \textit{Salmonella} $[2,2]$ and \textit{Campylobacter} $[3,3]$ from the lower tier $[4,6]$, whereas orderings among \textit{C. perfringens}, norovirus GI and norovirus GII remain unresolved. Surveillance comparisons should accompany reported hospitalization proportions with partial-identification bounds and joint rank confidence sets, and agencies prioritizing follow-up verification should target expected bound reduction per unit cost ($w_s U_i / (C_{js}\,\text{cost}_i)$) tailored to their operational cost structure.
